%----------Lista de Acronimos --------------------------------
\ifx\isenglish\undefined
\pretextualchapter{Lista de Siglas e Abreviaturas}
\else
\pretextualchapter{Acronyms}
\fi


%\renewcommand{\baselinestretch}{2} %espaçamento entre linhas (por definicçào, no espABNT é 2)
\begin{tabbing}
xxxxxxxxxxx \= xxxxxxxxxxxxx \kill
\textsc{API}            \> \textit{Application Programming Interface}\\


\textsc{BES}            \> Bacharelado em Engenharia de Software\\
\textsc{CA}            \> \textit{Certificate Authority} (autoridade certificadora)\\
\textsc{CI}            \> \textit{Continuous Integration} (integração contínua)\\
\textsc{CLI}            \> \textit{Command-Line Interface} (interface de linha de comando)\\
\textsc{HTTP}            \> \textit{Hypertext Transfer Protocol}\\
\textsc{HTTPS}            \> \textit{HTTP Secure}\\
\textsc{JSON}            \> \textit{JavaScript Object Notation}\\
\textsc{LLM}            \> \textit{Large Language Model} (modelo de linguagem de grande escala)\\
\textsc{MITM}            \> \textit{Man-in-the-Middle}\\
\textsc{MoSCoW}            \> \textit{Must, Should, Could, Won't} (técnica de priorização)\\
\textsc{PEM}            \> \textit{Privacy-Enhanced Mail} (formato de certificados e chaves)\\
\textsc{RF}            \> Requisito funcional\\
\textsc{RNF}            \> Requisito não funcional\\
\textsc{SGBD}            \> Sistema gerenciador de banco de dados\\
\textsc{TLS}            \> \textit{Transport Layer Security}\\
\textsc{TUI}            \> \textit{Terminal User Interface} (interface de terminal)\\
\textsc{UC}            \> Caso de uso\\
\textsc{YAML}            \> \textit{YAML Ain't Markup Language}

\end{tabbing}

%\addcontentsline{toc}{chapter}{Lista de Siglas e Abreviaturas}
